\documentclass[10pt,a4paper]{amsart}
\usepackage[margin=19mm]{geometry}
\usepackage{amsmath,amssymb,mathtools}
\usepackage[scr=rsfs,cal=euler]{mathalfa}
\usepackage{xfrac}
\usepackage[unicode,hidelinks,bookmarks=false]{hyperref}
\newcommand{\ie}{i.e.\ }
\newcommand{\cf}{cf.\ }
\newcommand{\resp}{respectively\ }
\newcommand{\Subsection}{\subsection}
\providecommand{\nobreakdash}{\nobreak\hbox{-}}
\setlength{\emergencystretch}{2em}
\multlinegap0pt
\usepackage{mathtools}
\usepackage{xcolor}
\usepackage{amssymb}
\usepackage{algorithm}\usepackage{algpseudocode}
\usepackage{tikz}
\usepackage[shortlabels]{enumitem}
\newtheorem{proposition}{Proposition}
\title{Dissipativity properties of a class of nonlinear time-delay systems via Bessel--Legendre inequalities}
\author{Ikram El Haskouki, Hannes Gernandt}
\date{Source: arXiv:2607.03808v1 (CC BY 4.0)}
\begin{document}
\maketitle

\noindent\emph{Source and licence.} Dissipativity properties of a class of nonlinear time-delay systems via Bessel--Legendre inequalities. Version arXiv:2607.03808v1 (CC BY 4.0), distributed by arXiv under the Creative Commons Attribution 4.0 International licence, http://creativecommons.org/licenses/by/4.0/, as recorded in the arXiv licence metadata for this exact version and shown on the abstract page https://arxiv.org/abs/2607.03808v1. Full licence notice: \texttt{licenses/arxiv-2607.03808-CC-BY-4.0.txt}.\par\smallskip

\noindent\textit{Editorial scope.} Counted unit: the article's proposition proposition (source0.tex lines 503--509). The article's complete original proof of it is delivered with it (source0.tex lines 510--519, one \texttt{proof} environment of 10 lines). No mathematics is written, repaired or re-derived: the statement and the proof are the article's own lines, in the article's own notation, and no retained line is substituted or re-flowed.  Retained uncounted as the source-local context the proof needs: lines 496--501.  Nothing was omitted from any retained span, so the package carries no omission record.  No retained line contains a citation, so the wrapper carries no bibliography and no bibliography entry.  No retained line contains a figure environment, an image-inclusion command or drawing code, so no image is delivered and the package declares no asset dependency.  Editorial formatting only, in addition: an AMS \texttt{amsart} preamble that restates the article's own declarations and macros used by the retained lines, namely the environments \texttt{proposition} on separate counters, together with the standard packages those lines need.  The retained environments print in this document as Proposition 1 in order of appearance, whereas the article prints the same items with the numbering of its own full document.  The complete native source of the article as distributed in the arXiv e-print for this exact version is bundled unchanged as \texttt{sources/204423/source0.tex}.
\begin{align}\label{eq:Interconnection}
\nonumber 
\dot x_j(t) &= f_j(x_j(t),x_j(t-\tau_{1,j}),\ldots, x_j(t-\tau_{q,j})) 
  + g_j(x_j(t))u_j(t), \quad t\geq 0,\\
y_j(t) &= h_j(x_j(t)),\quad \quad j=1,2,
\end{align}

\begin{proposition}
If the systems \eqref{eq:Interconnection} are impedance passive  with storage functions $\mathcal{S}_1$ and $\mathcal{S}_2$ then the negative feedback interconnection
\[
u_1=v-y_2, \quad u_2=y_1,
\]
with external input $v$ is impedance passive with storage function $\mathcal{S}=\mathcal{S}_1+\mathcal{S}_2$.  
\end{proposition}

\begin{proof}
The impedance passivity inequalities hold individually for arbitrary inputs which gives
\begin{align*}
\mathcal{S}(x_{t_1})-\mathcal{S}(x_{t_0})&=\mathcal{S}_1(x_{t_1}^1)-\mathcal{S}_1(x_{t_0}^1)+\mathcal{S}_2(x_{t_1}^2)-\mathcal{S}_2(x_{t_0}^2) \\ 
&\leq \int_{t_0}^{t_1}y_1(s)^\top u_1(s)\mathrm{d}s+\int_{t_0}^{t_1}y_2(s)^\top u_2(s)\mathrm{d}s\\
&=\int_{t_0}^{t_1}y_1(s)^\top (v-y_2)(s)\mathrm{d}s+\int_{t_0}^{t_1}y_2(s)^\top y_1(s)\mathrm{d}s\\
&=\int_{t_0}^{t_1}y_1(s)^\top v(s)\mathrm{d}s.
\end{align*}
Hence, the interconnected system is impedance passive with respect to the input $v$ and the output $y_1$.
\end{proof}

\end{document}
